\section{PROTOCOLO PRELIMINAR DE JUICIO DE EXPERTOS}
\label{anx:protocolo-expertos}

\begin{spacing}{1.5}
\subsection{Propósito y estado}

Borrador preliminar del protocolo del indicador F4 de la matriz de consistencia (Anexo~\ref{anx:matriz-consistencia}). Evalúa de forma estructurada las soluciones de cálculo térmico del artefacto frente a soluciones convencionales para cables enterrados. No se ha aplicado ni validado: el número de evaluadores, los pesos y el umbral de aceptación se fijan antes de la evaluación final.

\subsection{Objeto de comparación}

Los evaluadores comparan, con los mismos datos de entrada, el resultado del artefacto con la corriente admisible de la norma IEC 60287 \parencite{iec60287} y con una solución de elementos finitos verificada. Cada caso incluye geometría, materiales, pérdidas, contornos y resultado de ampacidad, sin indicar qué método produjo cada resultado.

\subsection{Evaluadores}

Panel de tres a cinco profesionales con experiencia en cables de potencia, ampacidad o modelado térmico. Cada evaluador declara por escrito sus vínculos con el estudio o con los métodos comparados, y no participa en los casos con conflicto de interés.

\subsection{Criterios y escala}

Cada caso se evalúa con cinco criterios, en una escala de 1 a 10:

\begin{enumerate}
  \item Exactitud y adecuación técnica: concordancia del resultado con la referencia y con el caso físico.
  \item Consistencia física: cumplimiento de la conservación de la energía, de la continuidad de temperatura y flujo en las interfaces, y de las condiciones de contorno.
  \item Utilidad para el cálculo de ampacidad: capacidad del resultado para sustentar la corriente admisible y el margen térmico.
  \item Trazabilidad: disponibilidad de datos, supuestos y versión suficientes para reconstruir el resultado.
  \item Claridad de límites: identificación explícita del dominio de validez y de las condiciones en las que el resultado no debe usarse.
\end{enumerate}

Anclajes: 1, el resultado contradice la referencia o la física sin explicación; 5, coincide en orden de magnitud, pero sus supuestos o límites no se documentan con claridad; 10, concuerda con la referencia dentro de la tolerancia del caso, respeta la física y declara sus límites con trazabilidad completa. Los valores intermedios se asignan por comparación con estos anclajes.

\subsection{Procedimiento}

Cada evaluador recibe los casos por separado y no conoce las puntuaciones de los demás; cuando el formato lo permite, el método se oculta. Registra una justificación escrita por criterio y caso. Si dos evaluadores difieren en más de tres puntos, revisan juntos la justificación; si la diferencia persiste, decide un tercero.

\subsection{Cuestionario por caso}

Antes de puntuar, el evaluador responde sí o no:

\begin{enumerate}
  \item ¿El caso declara geometría, materiales, pérdidas, contornos y versión?
  \item ¿El resultado reporta el error de $\Tmax$ frente a la referencia y el balance energético?
  \item ¿Los supuestos, la semilla y la configuración numérica están documentados?
  \item ¿Se declaran el dominio de validez y los usos no admitidos?
  \item ¿Cada puntuación tiene una justificación escrita?
\end{enumerate}

Un caso con respuestas negativas se reporta como inconcluso hasta completarse.

\subsection{Agregación y regla de aceptación}

La puntuación de cada criterio y caso es la mediana de los evaluadores. La concordancia se reporta con el coeficiente de correlación intraclase de dos vías, con acuerdo absoluto, cuyo tipo se fija antes de la aplicación. El umbral mínimo por criterio y caso se establece antes de la evaluación final; los casos sin consenso se reportan como inconclusos.

\subsection{Relación con las métricas numéricas}

El juicio de expertos aporta una evaluación de utilidad y de adecuación profesional. Las métricas numéricas (error de $\Tmax$, NRMSE, desequilibrio energético y $\Delta I$) sustentan la verificación. Ninguna fuente reemplaza a la otra. Una puntuación alta no compensa un incumplimiento numérico, y un resultado numérico correcto no basta si su utilidad o sus límites no están documentados.

\subsection{Limitaciones del protocolo}

El juicio depende de la experiencia de los evaluadores y de la calidad del material entregado, y su aplicación exige tiempo y disponibilidad. Por ello, el protocolo se revisa antes de la aplicación, y sus resultados se reportan como evidencia complementaria.
\end{spacing}
